\documentclass[../../main]{subfiles}

\begin{document}

\section{Funções Inversas}
\label{sec:matematica:funcoes:funcoes-inversas}

\subsection{Funções Invertíveis}

\begin{defn}[Função Invertível]
  \label{def:matematica:funcoes:funcoes-inversas:invertivel}

  Seja

  \[
    f:A\to B.
  \]

  Diz-se que \(f\) é invertível quando existe uma função

  \[
    g:B\to A
  \]

  tal que

  \[
    g\circ f
    =
    \operatorname{id}_A
  \]

  e

  \[
    f\circ g
    =
    \operatorname{id}_B.
  \]

\end{defn}

\begin{defn}[Função Inversa]
  \label{def:matematica:funcoes:funcoes-inversas:funcao-inversa}

  Seja

  \[
    f:A\to B.
  \]

  Se \(f\) é invertível, a única função

  \[
    g:B\to A
  \]

  que satisfaz

  \[
    g\circ f
    =
    \operatorname{id}_A
  \]

  e

  \[
    f\circ g
    =
    \operatorname{id}_B
  \]

  é denominada função inversa de \(f\).

  Escreve-se

  \[
    g=f^{-1}.
  \]

\end{defn}

\begin{prop}[Unicidade da Inversa]
  \label{prop:matematica:funcoes:funcoes-inversas:unicidade}

  Uma função invertível possui uma única inversa.

\end{prop}

\begin{proof}

  Sejam

  \[
    g,h:B\to A
  \]

  inversas de \(f\).

  Então

  \[
    g
    =
    g\circ\operatorname{id}_B.
  \]

  Como

  \[
    \operatorname{id}_B
    =
    f\circ h,
  \]

  segue que

  \[
    g
    =
    g\circ(f\circ h).
  \]

  Pela associatividade,

  \[
    g
    =
    (g\circ f)\circ h.
  \]

  Como

  \[
    g\circ f
    =
    \operatorname{id}_A,
  \]

  temos

  \[
    g
    =
    \operatorname{id}_A\circ h.
  \]

  Logo

  \[
    g=h.
  \]

\end{proof}

\subsection{Invertibilidade e Bijetividade}

\begin{prop}
  \label{prop:matematica:funcoes:funcoes-inversas:invertivel-implica-injetiva}

  Toda função invertível é injetiva.

\end{prop}

\begin{proof}

  Seja

  \[
    f(x)=f(y).
  \]

  Aplicando \(f^{-1}\) em ambos os lados,

  \[
    f^{-1}(f(x))
    =
    f^{-1}(f(y)).
  \]

  Logo

  \[
    x=y.
  \]

\end{proof}

\begin{prop}
  \label{prop:matematica:funcoes:funcoes-inversas:invertivel-implica-sobrejetiva}

  Toda função invertível é sobrejetiva.

\end{prop}

\begin{proof}

  Seja

  \[
    y\in B.
  \]

  Defina

  \[
    x=f^{-1}(y).
  \]

  Então

  \[
    f(x)
    =
    f(f^{-1}(y))
    =
    y.
  \]

  Logo \(f\) é sobrejetiva.

\end{proof}

\begin{prop}
  \label{prop:matematica:funcoes:funcoes-inversas:invertivel-implica-bijetiva}

  Toda função invertível é bijetiva.

\end{prop}

\begin{proof}

  Segue das duas proposições anteriores.

\end{proof}

\begin{prop}
  \label{prop:matematica:funcoes:funcoes-inversas:bijetiva-implica-invertivel}

  Toda função bijetiva é invertível.

\end{prop}

\begin{proof}

  Seja

  \[
    f:A\to B
  \]

  bijetiva.

  Pela sobrejetividade, para cada

  \[
    y\in B
  \]

  existe ao menos um

  \[
    x\in A
  \]

  tal que

  \[
    f(x)=y.
  \]

  Pela injetividade, esse elemento é único.

  Definimos

  \[
    f^{-1}(y)=x.
  \]

  Assim obtemos uma função

  \[
    f^{-1}:B\to A.
  \]

  Pela construção,

  \[
    f^{-1}(f(x))
    =
    x
  \]

  e

  \[
    f(f^{-1}(y))
    =
    y.
  \]

  Portanto

  \[
    f^{-1}\circ f
    =
    \operatorname{id}_A
  \]

  e

  \[
    f\circ f^{-1}
    =
    \operatorname{id}_B.
  \]

\end{proof}

\begin{thm}[Caracterização Fundamental]
  \label{thm:matematica:funcoes:funcoes-inversas:caracterizacao}

  Seja

  \[
    f:A\to B.
  \]

  Então

  \[
    f
    \text{ é invertível}
    \iff
    f
    \text{ é bijetiva}.
  \]

\end{thm}

\begin{proof}

  Combinação das proposições anteriores.

\end{proof}

\subsection{Propriedades}

\begin{prop}
  \label{prop:matematica:funcoes:funcoes-inversas:inversa-da-inversa}

  Se \(f\) é invertível, então

  \[
    (f^{-1})^{-1}
    =
    f.
  \]

\end{prop}

\begin{proof}

  Como

  \[
    f^{-1}\circ f
    =
    \operatorname{id}_A
  \]

  e

  \[
    f\circ f^{-1}
    =
    \operatorname{id}_B,
  \]

  segue que \(f\) é inversa de \(f^{-1}\).

  Pela unicidade da inversa,

  \[
    (f^{-1})^{-1}
    =
    f.
  \]

\end{proof}

\begin{prop}[Inversa da Composição]
  \label{prop:matematica:funcoes:funcoes-inversas:inversa-composicao}

  Se

  \[
    f:A\to B
  \]

  e

  \[
    g:B\to C
  \]

  são invertíveis, então

  \[
    (g\circ f)^{-1}
    =
    f^{-1}\circ g^{-1}.
  \]

\end{prop}

\begin{proof}

  Basta verificar que

  \[
    (f^{-1}\circ g^{-1})
    \circ
    (g\circ f)
    =
    \operatorname{id}_A
  \]

  e

  \[
    (g\circ f)
    \circ
    (f^{-1}\circ g^{-1})
    =
    \operatorname{id}_C.
  \]

  A demonstração segue da associatividade da composição.

\end{proof}

\begin{prop}
  \label{prop:matematica:funcoes:funcoes-inversas:potencia}

  Se \(f:A\to A\) é invertível, então

  \[
    (f^n)^{-1}
    =
    (f^{-1})^n
  \]

  para todo

  \[
    n\in\mathbb N.
  \]

\end{prop}

\begin{proof}

  Por indução em \(n\).

\end{proof}

\subsection{Permutações}

\begin{prop}
  \label{prop:matematica:funcoes:funcoes-inversas:permutacao}

  Toda permutação possui inversa, que também é uma permutação.

\end{prop}

\begin{proof}

  Uma permutação é uma bijeção

  \[
    f:A\to A.
  \]

  Pelo teorema anterior, \(f\) é invertível.

  Sua inversa também é uma bijeção de \(A\) em \(A\).

\end{proof}

\end{document}
